\chapter{1922:Francis W. Aston}
\label{chap:chemistry1922}

The Nobel Prize in Chemistry for the year 1922 was awarded to Francis W. Aston.

\textbf{Motivation:}
for his discovery, by means of his mass spectrograph, of isotopes, in a large number of non-radioactive elements, and for his enunciation of the whole-number rule

\section{Francis W. Aston}

\subsection*{Early Life and Education}

Francis W. Aston was born on 1877-09-01 in Harborne, United Kingdom.

\subsection*{Career and Major Research}

Aston worked at the Cavendish Laboratory in Cambridge. He built the mass spectrograph, an instrument that separated atomic nuclei according to their mass, and used it to measure the exact masses of the isotopes of many elements. His work established the whole-number rule, according to which isotope masses are nearly integers relative to hydrogen.

\subsection*{Nobel Prize-winning Work}

The work for which Francis W. Aston was awarded the Nobel Prize in Chemistry in 1922 was described by the Nobel Committee as: "for his discovery, by means of his mass spectrograph, of isotopes, in a large number of non-radioactive elements, and for his enunciation of the whole-number rule".

This research fundamentally advanced the field by making groundbreaking contributions that transformed chemical science.

\subsection*{Significance and Impact}

The impact of Francis W. Aston's work extends far beyond the specific achievement recognized by the Nobel Prize.
These contributions continue to influence chemical research and education today.

\subsection*{Later Career and Honors}

Francis W. Aston continued his scientific work until 1945-11-20, passing away in Cambridge, United Kingdom.
He received numerous honors including membership in major scientific academies and societies.

\subsection*{Legacy}

The legacy of Francis W. Aston endures through the lasting impact of his scientific contributions.
Francis W. Aston's work continues to be cited and built upon by researchers across the chemical sciences.

\subsection*{Selected Publications}

Francis W. Aston's major publications include numerous papers in leading journals of the era, detailing the experimental and theoretical work summarized above.

\clearpage

\section*{Chapter Summary}

Francis W. Aston received the prize for his discovery, by means of the mass spectrograph, of isotopes in a large number of non-radioactive elements.
